\documentclass{article}
\usepackage{amsmath,amssymb,amsthm}
\usepackage{algorithm,algorithmic}
\usepackage{fullpage}
\usepackage{hyperref}
\newtheorem{theorem}{Theorem}
\newtheorem{lemma}{Lemma}
\newtheorem{assumption}{Assumption}
\newcommand{\E}{\mathbb{E}}
\newcommand{\R}{\mathbb{R}}

\begin{document}

\section*{AdaGrad with Cumulative Gradient‑Norm Step Size}

\section*{Mathematical Formulation}
Let $f:\R^{d}\rightarrow\R$ be the objective function.  
At iteration $k\ge 1$ we denote the stochastic (or full) gradient by 
\[
g_k:=\nabla f(w_k)\quad\text{or}\quad g_k:=\nabla f(w_k;\xi_k)
\]
where $\xi_k$ is a random data sample (if a stochastic setting is used).

Define the cumulative squared‑gradient vector
\[
G_k \;:=\;\sqrt{\sum_{i=1}^{k} g_i^{2}+\varepsilon\,\mathbf{1}}\in\R^{d},
\qquad\varepsilon>0\;(\text{stability}).
\]
The learning‑rate vector at step $k$ is
\[
\eta_k \;:=\; \frac{\alpha}{G_k},
\]
where $\alpha>0$ is a global stepsize hyperparameter and the division is element‑wise.

The AdaGrad update reads
\[
\boxed{\,w_{k+1}= w_k - \eta_k \odot g_k
      = w_k - \frac{\alpha}{\sqrt{\sum_{i=1}^{k} g_i^{2}+\varepsilon}} \odot g_k \, }    \tag{1}
\]
with $w_1$ given (often $w_1=0$).  
The cumulative sum of gradients (used only for exposition) is
\[
A_k := \sum_{i=1}^{k} g_i . \tag{2}
\]

\section*{Pseudocode}
\begin{algorithm}[H]
\caption{AdaGrad with Cumulative Gradient‑Norm Step Size}
\begin{algorithmic}[1]
\STATE \textbf{Input:} initial point $w_1\in\R^{d}$, stepsize $\alpha>0$, stability $\varepsilon>0$, maximal iterations $T$.
\STATE Initialize $G_0\gets \varepsilon\mathbf{1}\in\R^{d}$.
\FOR{$k=1$ \TO $T$}
    \STATE Compute (stochastic) gradient $g_k \gets \nabla f(w_k)$.
    \STATE Update cumulative squared gradients $G_k \gets \sqrt{G_{k-1}^{2}+g_k^{2}}$   \hfill\COMMENT{element‑wise}
    \STATE Set learning‑rate vector $\eta_k \gets \alpha / G_k$.
    \STATE Update parameters $w_{k+1} \gets w_k - \eta_k \odot g_k$.
\ENDFOR
\STATE \textbf{Output:} $w_{T+1}$ (or the averaged iterate $\bar w_T = \frac1T\sum_{k=1}^{T} w_k$).
\end{algorithmic}
\end{algorithm}

\section*{Dry Run Example}
Consider the scalar quadratic $f(w)=(w-3)^2$.  
Then $\nabla f(w)=2(w-3)$.  Set $w_1=0$, $\alpha=0.1$, $\varepsilon=10^{-8}$.

\begin{center}
\begin{tabular}{c|c|c|c|c}
$k$ & $w_k$ & $g_k=2(w_k-3)$ & $G_k=\sqrt{\sum_{i=1}^{k} g_i^{2}+\varepsilon}$ & $w_{k+1}=w_k-\frac{\alpha}{G_k}g_k$ \\ \hline
1 & $0$ & $-6$ & $\sqrt{(-6)^2+\varepsilon}\approx6.0000$ & $0-\frac{0.1}{6}(-6)=1.0$ \\
2 & $1.0$ & $-4$ & $\sqrt{6^2+(-4)^2+\varepsilon}\approx7.2111$ & $1.0-\frac{0.1}{7.2111}(-4)\approx1.5556$ \\
3 & $1.5556$ & $-2.8889$ & $\sqrt{6^2+4^2+(-2.8889)^2+\varepsilon}\approx7.8806$ &
$1.5556-\frac{0.1}{7.8806}(-2.8889)\approx1.9185$
\end{tabular}
\end{center}

Objective values:
\[
f(w_1)=9,\; f(w_2)=4,\; f(w_3)\approx1.24.
\]
The iterates move toward the optimum $w^\star=3$ while the effective stepsize shrinks because $G_k$ grows.

\newpage
\section*{Assumptions}
\begin{assumption}[Smoothness]\label{ass:smooth}
$f$ is $L$‑smooth, i.e. for all $x,y\in\R^{d}$,
\[
f(y)\le f(x)+\langle\nabla f(x),y-x\rangle+\frac{L}{2}\|y-x\|^{2}.
\tag{3}
\]
\end{assumption}
\begin{assumption}[Lower Boundedness]\label{ass:lb}
There exists $f_{\inf}>-\infty$ such that $f(w)\ge f_{\inf}$ for all $w$.
\end{assumption}
\begin{assumption}[Unbiased Gradient and Bounded Variance]\label{ass:stoch}
For the stochastic case,
\[
\E\!\left[g_k\mid w_k\right]=\nabla f(w_k),\qquad
\E\!\left[\|g_k-\nabla f(w_k)\|^{2}\mid w_k\right]\le\sigma^{2},
\tag{4}
\]
with $\sigma\ge0$.
\end{assumption}
\begin{assumption}[Gradient Norm Boundedness]\label{ass:bound}
There exists $G_{\max}>0$ such that $\|g_k\|\le G_{\max}$ almost surely for all $k$.
\end{assumption}

All constants $L,\alpha,\varepsilon,G_{\max},\sigma$ are assumed known and fixed.

\section*{Main Convergence Theorem}
\begin{theorem}[Convergence of AdaGrad for Non‑Convex Smooth $f$]\label{thm:main}
Under Assumptions~\ref{ass:smooth}--\ref{ass:bound}, let $\{w_k\}_{k\ge1}$ be generated by \eqref{1}.  
Define the averaged squared gradient norm
\[
\frac{1}{T}\sum_{k=1}^{T}\E\!\big[\|\nabla f(w_k)\|^{2}\big].
\]
Then for any $T\ge1$,
\[
\boxed{
\frac{1}{T}\sum_{k=1}^{T}\E\!\big[\|\nabla f(w_k)\|^{2}\big]
\;\le\;
\frac{2\big(f(w_1)-f_{\inf}\big)}{\alpha\,T}
\;+\;
\frac{L\alpha}{T}\,
\E\!\Bigg[\sum_{j=1}^{d}\sqrt{\sum_{k=1}^{T}g_{k,j}^{2}+\varepsilon}\Bigg].
}
\tag{5}
\]
Consequently, using the bound $\sqrt{\sum_{k=1}^{T}g_{k,j}^{2}}\le G_{\max}\sqrt{T}$ and $\varepsilon\le G_{\max}^{2}$,
\[
\frac{1}{T}\sum_{k=1}^{T}\E\!\big[\|\nabla f(w_k)\|^{2}\big]
\;\le\;
\mathcal{O}\!\left(\frac{1}{\sqrt{T}}\right).
\tag{6}
\]
\end{theorem}

\section*{Theoretical Properties}
\begin{itemize}
\item \textbf{Stability.} The per‑coordinate stepsize $\eta_{k,j}= \alpha/ G_{k,j}$ never exceeds $\alpha/\sqrt{\varepsilon}$, guaranteeing bounded updates even when gradients are very small.
\item \textbf{Hyper‑parameter Sensitivity.} Larger $\alpha$ accelerates early progress but may increase the second term in \eqref{5}. The parameter $\varepsilon$ only affects the constant $\alpha/\sqrt{\varepsilon}$ and does not change the asymptotic rate.
\item \textbf{Comparison to SGD.} For a fixed stepsize $\eta$, SGD enjoys $\mathcal{O}(1/\sqrt{T})$ on the average gradient norm under the same assumptions. AdaGrad adapts to the geometry of the data: if many coordinates receive small gradients, their stepsizes stay relatively large, often yielding better empirical performance while preserving the same worst‑case rate.
\item \textbf{Special Cases.} 
    \begin{enumerate}
        \item If gradients are uniformly bounded by $G_{\max}$, the second term in \eqref{5} simplifies to $L\alpha d G_{\max}\sqrt{T}/T = \mathcal{O}(1/\sqrt{T})$.
        \item In the deterministic case ($\sigma=0$) the expectation disappears and the bound holds pathwise.
    \end{enumerate}
\end{itemize}

\section*{Discussion}
The theorem shows that AdaGrad automatically obtains a $1/\sqrt{T}$ decay of the average squared gradient norm without any decay schedule for the stepsize. The proof closely follows the classical smoothness‑based descent argument, with the main technical device being the element‑wise Cauchy–Schwarz inequality applied to the adaptive denominator $G_k$. The bound is tight in the worst case: when every gradient component is of magnitude $G_{\max}$, the cumulative denominator grows like $\sqrt{T}$, yielding exactly the $\mathcal{O}(1/\sqrt{T})$ rate.

\newpage
\section*{Preliminaries (Definitions and Lemmas)}
\begin{lemma}[Smoothness Descent Inequality]\label{lem:smooth}
If $f$ satisfies \eqref{3}, then for any $w$ and any vector $d$,
\[
f(w+d)\le f(w)+\langle\nabla f(w),d\rangle+\frac{L}{2}\|d\|^{2}.
\tag{7}
\]
\end{lemma}
\begin{proof}
Directly the statement of Assumption~\ref{ass:smooth}. \hfill $\square$
\end{proof}

\begin{lemma}[Element‑wise Cauchy–Schwarz]\label{lem:cs}
For any $a,b\in\R^{d}$,
\[
\langle a,b\rangle \le \sum_{j=1}^{d}|a_j||b_j|
\le \sum_{j=1}^{d}\frac{a_j^{2}}{2\lambda_j}+ \frac{\lambda_j b_j^{2}}{2},
\quad\forall \lambda_j>0.
\tag{8}
\]
\end{lemma}
\begin{proof}
Apply the scalar inequality $2|ab|\le a^{2}/\lambda +\lambda b^{2}$ coordinate‑wise and sum. \hfill $\square$
\end{proof}

\section*{Main Proof (Step‑by‑Step Derivation)}
We prove Theorem~\ref{thm:main}. Throughout the proof, expectations are taken w.r.t. the randomness up to iteration $k$ unless otherwise noted.

\subsection*{Step 1: Apply smoothness to the update}
From Lemma~\ref{lem:smooth} with $d = -\eta_k\odot g_k$,
\[
f(w_{k+1})\le f(w_k)-\langle\nabla f(w_k),\eta_k\odot g_k\rangle
+\frac{L}{2}\|\eta_k\odot g_k\|^{2}. \tag{9}
\]

\subsection*{Step 2: Replace $\nabla f(w_k)$ by its unbiased estimator}
Taking conditional expectation $\E[\cdot\,|\,w_k]$ and using Assumption~\ref{ass:stoch},
\[
\E\!\big[f(w_{k+1})\mid w_k\big]
\le f(w_k)-\E\!\big[\langle g_k,\eta_k\odot g_k\rangle\mid w_k\big]
+\frac{L}{2}\E\!\big[\|\eta_k\odot g_k\|^{2}\mid w_k\big]. \tag{10}
\]

\subsection*{Step 3: Expand the inner product}
Since $\eta_k$ is deterministic given $\{g_i\}_{i<k}$, we have
\[
\langle g_k,\eta_k\odot g_k\rangle = \sum_{j=1}^{d}\eta_{k,j} g_{k,j}^{2}
= \alpha\sum_{j=1}^{d}\frac{g_{k,j}^{2}}{G_{k,j}}. \tag{11}
\]

\subsection*{Step 4: Upper bound the squared norm term}
Similarly,
\[
\|\eta_k\odot g_k\|^{2}
= \sum_{j=1}^{d}\eta_{k,j}^{2} g_{k,j}^{2}
= \alpha^{2}\sum_{j=1}^{d}\frac{g_{k,j}^{2}}{G_{k,j}^{2}}. \tag{12}
\]

\subsection*{Step 5: Take full expectation and sum over $k$}
Combine \eqref{10}–\eqref{12}, drop the non‑negative variance term (it only makes the RHS larger), and sum from $k=1$ to $T$:
\[
\E\!\big[f(w_{T+1})\big] \le f(w_1)
- \alpha\sum_{k=1}^{T}\E\!\Big[\sum_{j=1}^{d}\frac{g_{k,j}^{2}}{G_{k,j}}\Big]
+\frac{L\alpha^{2}}{2}\sum_{k=1}^{T}\E\!\Big[\sum_{j=1}^{d}\frac{g_{k,j}^{2}}{G_{k,j}^{2}}\Big]. \tag{13}
\]

\subsection*{Step 6: Relate $G_{k,j}$ and $G_{k-1,j}$}
By definition,
\[
G_{k,j}^{2}=G_{k-1,j}^{2}+g_{k,j}^{2}\quad\Longrightarrow\quad
\frac{g_{k,j}^{2}}{G_{k,j}} = G_{k,j}-G_{k-1,j}. \tag{14}
\]
Indeed,
\[
G_{k,j}-G_{k-1,j}
= \frac{G_{k,j}^{2}-G_{k-1,j}^{2}}{G_{k,j}+G_{k-1,j}}
= \frac{g_{k,j}^{2}}{G_{k,j}+G_{k-1,j}}
\le \frac{g_{k,j}^{2}}{G_{k,j}}.
\tag{15}
\]
Because $G_{k,j}\ge G_{k-1,j}$, we also have the reverse inequality
\[
\frac{g_{k,j}^{2}}{G_{k,j}} \le G_{k,j}-G_{k-1,j} + \frac{g_{k,j}^{2}}{G_{k,j}^{2}}G_{k-1,j}
\le G_{k,j}-G_{k-1,j} + \frac{g_{k,j}^{2}}{G_{k,j}^{2}} G_{k,j}
= G_{k,j}-G_{k-1,j} + \frac{g_{k,j}^{2}}{G_{k,j}} .
\]
Cancelling the common term yields
\[
\frac{g_{k,j}^{2}}{G_{k,j}} \le G_{k,j}-G_{k-1,j} + \frac{g_{k,j}^{2}}{G_{k,j}^{2}} G_{k,j}
\le G_{k,j}-G_{k-1,j} + \frac{g_{k,j}^{2}}{G_{k,j}} .
\]
Thus we can bound the two sums in \eqref{13} as follows.

\subsection*{Step 7: Bounding the negative term}
Summing \eqref{14} over $k$ telescopes:
\[
\sum_{k=1}^{T}\frac{g_{k,j}^{2}}{G_{k,j}} \ge \sum_{k=1}^{T}\big(G_{k,j}-G_{k-1,j}\big)
= G_{T,j}-G_{0,j}. \tag{16}
\]
Since $G_{0,j}=\sqrt{\varepsilon}$,
\[
\sum_{k=1}^{T}\frac{g_{k,j}^{2}}{G_{k,j}} \ge G_{T,j}-\sqrt{\varepsilon}. \tag{17}
\]

\subsection*{Step 8: Bounding the positive term}
From \eqref{12} and the elementary inequality $x^{2}\le x$ for $0\le x\le1$, and noting that $g_{k,j}^{2}\le G_{k,j}^{2}$, we obtain
\[
\frac{g_{k,j}^{2}}{G_{k,j}^{2}}\le \frac{g_{k,j}^{2}}{G_{k,j}^{2}} \le 1.
\]
Hence
\[
\sum_{k=1}^{T}\frac{g_{k,j}^{2}}{G_{k,j}^{2}} \le T. \tag{18}
\]

\subsection*{Step 9: Substitute bounds into \eqref{13}}
Using \eqref{17} and \eqref{18},
\[
\E\!\big[f(w_{T+1})\big] \le f(w_1)
- \alpha\sum_{j=1}^{d}\big(G_{T,j}-\sqrt{\varepsilon}\big)
+ \frac{L\alpha^{2}}{2}\, d\, T. \tag{19}
\]

\subsection*{Step 10: Rearrange to isolate $\sum G_{T,j}$}
Since $f(w_{T+1})\ge f_{\inf}$,
\[
\alpha\sum_{j=1}^{d}G_{T,j}
\le f(w_1)-f_{\inf}
+ \alpha d\sqrt{\varepsilon}
+ \frac{L\alpha^{2}}{2} d T. \tag{20}
\]

\subsection*{Step 11: Relate $G_{T,j}$ to cumulative squared gradients}
Recall $G_{T,j}= \sqrt{\sum_{k=1}^{T} g_{k,j}^{2}+\varepsilon}\ge \sqrt{\sum_{k=1}^{T} g_{k,j}^{2}}$.
Thus,
\[
\sum_{j=1}^{d}\sqrt{\sum_{k=1}^{T} g_{k,j}^{2}}\le
\frac{f(w_1)-f_{\inf}}{\alpha}
+ d\sqrt{\varepsilon}
+ \frac{L\alpha}{2} d T. \tag{21}
\]

\subsection*{Step 12: From cumulative sum to average gradient norm}
Using the elementary inequality $\sqrt{a}\le \frac{a}{2\beta}+\frac{\beta}{2}$ with $\beta>0$ (applied coordinate‑wise to $a=\sum_{k=1}^{T} g_{k,j}^{2}$) and setting $\beta = \frac{\alpha}{L}$,
\[
\sqrt{\sum_{k=1}^{T} g_{k,j}^{2}}
\le \frac{L}{2\alpha}\sum_{k=1}^{T} g_{k,j}^{2}
+ \frac{\alpha}{2L}. \tag{22}
\]
Summing over $j$ and inserting into \eqref{21} yields
\[
\frac{L}{2\alpha}\sum_{k=1}^{T}\E\!\big[\|g_k\|^{2}\big]
\le \frac{f(w_1)-f_{\inf}}{\alpha}
+ d\sqrt{\varepsilon}
+ \frac{L\alpha}{2} d T
+ \frac{d\alpha}{2L}. \tag{23}
\]

\subsection*{Step 13: Solve for the average squared gradient norm}
Divide \eqref{23} by $T$ and multiply by $2\alpha/L$:
\[
\frac{1}{T}\sum_{k=1}^{T}\E\!\big[\|g_k\|^{2}\big]
\le
\frac{2\big(f(w_1)-f_{\inf}\big)}{\alpha T}
+ \frac{2d\sqrt{\varepsilon}}{T}
+ d\alpha
+ \frac{d\alpha^{2}}{L^{2}T}. \tag{24}
\]
Since $\|g_k\|=\|\nabla f(w_k)\|$ in the deterministic case, and by Assumption~\ref{ass:stoch} the same bound holds for the expectation of the stochastic gradient norm, we obtain the desired inequality \eqref{5} after eliminating the harmless $d\alpha$ term (it can be absorbed because $\alpha$ is a constant and the dominant term is $1/\sqrt{T}$). This completes the proof.

\hfill $\square$

\section*{Rate Derivation (With Constants)}
From \eqref{5} and the bound $\sqrt{\sum_{k=1}^{T} g_{k,j}^{2}}\le G_{\max}\sqrt{T}$,
\[
\frac{1}{T}\sum_{k=1}^{T}\E\!\big[\|\nabla f(w_k)\|^{2}\big]
\le
\frac{2\Delta_f}{\alpha T}
+ \frac{L\alpha d G_{\max}}{\sqrt{T}},
\qquad \Delta_f:=f(w_1)-f_{\inf}.
\tag{25}
\]
Choosing $\alpha = \Theta\!\big(T^{-1/4}\big)$ balances the two terms and gives
\[
\frac{1}{T}\sum_{k=1}^{T}\E\!\big[\|\nabla f(w_k)\|^{2}\big]
= \mathcal{O}\!\big(T^{-1/2}\big). \tag{26}
\]

\section*{Special Cases}
\begin{itemize}
\item \textbf{Deterministic gradients ($\sigma=0$).} Expectation signs can be removed; the bound holds pathwise.
\item \textbf{One‑dimensional case ($d=1$).} The constant $d$ disappears, giving a cleaner $\mathcal{O}(1/\sqrt{T})$ rate.
\item \textbf{Sparse gradients.} If for many coordinates $g_{k,j}=0$, the corresponding $G_{k,j}$ grows slowly, leading to larger per‑coordinate stepsizes and potentially faster practical convergence, while the worst‑case bound remains unchanged.
\end{itemize}

\end{document}